\chapter{Системы контроля версий и процессы разработки}
\label{ch:12}

Предыдущие главы учили писать код, и всё это время программа была текстом, который лежит в папке на вашем компьютере. Эта глава смотрит на неё иначе~--- как на продукт, который делает команда и который нужно доводить до пользователя без потерь. Вы разберёте три связанные вещи: Git~--- распределённую систему контроля версий, ставшую стандартом отрасли; CI/CD~--- автоматические конвейеры, проверяющие код после каждого изменения; Agile и Scrum~--- правила, по которым команда планирует работу. Темы неразделимы: Git даёт технику совместной работы, CI/CD~--- автоматизацию проверок, Scrum~--- организацию людей вокруг всего этого.

\section{Зачем нужна система контроля версий}

\begin{sloppypar}
Вспомните, как пишется курсовая. Сначала появляется файл \texttt{kursovaya.docx}, потом \texttt{kursovaya\_v2.docx}, затем \texttt{kursovaya\_final.docx} и, разумеется, \texttt{kursovaya\_final\_final.docx}. А когда руководитель просит вернуть абзац, удалённый три дня назад, выясняется, что его больше нет нигде. Теперь представьте, что так работают пять человек над проектом из трёхсот файлов.
\end{sloppypar}

\term{Система контроля версий} (\eng{Version Control System, VCS})~--- программа, которая хранит полную историю изменений набора файлов, позволяет вернуться к любому прошлому состоянию, сравнить состояния и объединить работу нескольких человек. Отсюда всё остальное: видно, кто и зачем изменил каждую строку, а поверх истории строятся код-ревью, автоматические проверки и выпуск релизов.

Систем два семейства. \term{Централизованные} (Subversion, Perforce) хранят историю на единственном сервере, а у разработчика лежит только рабочая копия последней версии. \term{Распределённые} (Git, Mercurial) дают каждому участнику полную копию репозитория со всей историей. Разница как между библиотекой с единственным экземпляром книги и домашней копией всей библиотеки у каждого читателя: сгорела библиотека~--- книги больше нет. Следствия собраны в табл.~\ref{tab:12-1}.

\begin{table}[htbp]
\caption{Централизованные и распределённые системы контроля версий}
\label{tab:12-1}
\centering
\begin{tabular}{L{3.6cm}L{5.4cm}L{5.4cm}}
\toprule
Критерий & Централизованная & Распределённая (Git) \\
\midrule
Хранение истории & только на сервере & полностью у каждого участника \\
Работа без сети & почти невозможна & полноценная, кроме обмена \\
Скорость операций & зависит от сети & локальные операции мгновенны \\
Ветвление & переключение и слияние идут через сервер & всё локально: указатель на коммит \\
Отказ сервера & работа останавливается & работа продолжается локально \\
\bottomrule
\end{tabular}
\end{table}

Именно дешёвое ветвление сделало Git стандартом: всё происходит на вашем диске без обращения к серверу, поэтому ветку заводят под каждую задачу. Git создал Линус Торвальдс в 2005~году для работы над ядром Linux; после установки его настраивают один раз (листинг~\ref{lst:12-1}): имя и почта попадут в каждый ваш коммит.

\begin{listing}[H]
\begin{minted}{bash}
git config --global user.name "Ivan Ivanov"
git config --global user.email "ivan@example.com"
git config --global init.defaultBranch main
git config --list                 # проверить результат
\end{minted}
\caption{Первичная настройка Git}
\label{lst:12-1}
\end{listing}

Отдельная деталь~--- окончания строк: Windows завершает строку парой символов CRLF, а Linux и macOS~--- одним символом LF. Если этого не настроить, файл после переноса между системами превращается в разницу, где изменённой помечена буквально каждая строка. Лечит настройка \texttt{core.autocrlf} (\texttt{true} на Windows, \texttt{input} на Linux и macOS), а надёжнее~--- файл \texttt{.gitattributes} со строкой \texttt{* text=auto} в самом репозитории: тогда правило действует для всех.

\section{Модель данных Git}

Большинство «магических» проблем Git перестают быть магическими, как только вы представляете его модель данных. Старые системы хранят историю как список правок; Git устроен иначе: каждый коммит~--- это \term{снимок} (\eng{snapshot}) всего проекта целиком. Обычная система ведёт дневник ремонта, а Git каждый день фотографирует всю квартиру; если комната не менялась, снимок ссылается на прежнюю её фотографию, поэтому репозиторий не раздувается. Восстановить версию~--- значит взять готовый снимок, а не применить четыреста правок подряд; отсюда и скорость переключения между ветками.

База объектов состоит из четырёх типов записей. \term{Blob} хранит содержимое одного файла~--- без имени и прав доступа. \term{Tree} описывает каталог: какому имени какой blob или tree соответствует. \term{Commit} ссылается на корневой tree и на родительские коммиты, хранит автора, дату и сообщение. \term{Tag}~--- аннотированная метка со ссылкой на объект. Связи объектов внутри одного коммита показаны на рис.~\ref{fig:12-1}: имя файла записано не в blob, а в дереве, поэтому два одинаковых по содержимому файла с разными именами дают один blob и две записи в дереве.

\begin{figure}[htbp]\centering
\begin{tikzpicture}[font=\footnotesize,
  ob/.style={draw,rounded corners=2pt,align=center,inner sep=4pt,
    font=\footnotesize}]
\node[ob] (c) {\textbf{commit} \texttt{a3f1c9e}\\[1pt]автор, дата, сообщение};
\node[ob,right=16mm of c] (t1) {\textbf{tree} \texttt{92b8b69}\\[1pt]корень
  проекта};
\node[ob,right=16mm of t1] (b1) {\textbf{blob} \texttt{ce01362}\\[1pt]%
  \texttt{README.md}};
\node[ob,below=7mm of b1] (t2) {\textbf{tree} \texttt{4d7a2b1}\\[1pt]%
  \texttt{src/}};
\node[ob,below=7mm of t2] (b2) {\textbf{blob} \texttt{8ab3f01}\\[1pt]%
  \texttt{Main.java}};
\node[ob,below=13mm of c] (p) {\textbf{commit} \texttt{8f2c1a4}\\[1pt]%
  родительский коммит};
\draw[flowarrow] (c) -- (t1);
\draw[flowarrow] (t1) -- (b1);
\draw[flowarrow] (t1) -- (t2);
\draw[flowarrow] (t2) -- (b2);
\draw[flowarrow] (c) -- (p);
\end{tikzpicture}
\caption{Объекты, из которых состоит один коммит}\label{fig:12-1}
\end{figure}

Имя каждого объекта~--- криптографический хеш его содержимого, сорок шестнадцатеричных символов алгоритма SHA-1, работающий как отпечаток пальца. Хеш коммита считается в том числе от хеша родителя, поэтому изменить старый коммит незаметно нельзя~--- поменяются хеши всех последующих. А один и тот же коммит на вашей машине и на сервере имеет один идентификатор, так что централизованная нумерация версий не нужна.

Поверх объектов Git держит \term{ссылки} (\eng{refs})~--- текстовые файлы с хешем коммита внутри. \term{Ветка}~--- ссылка, которая сдвигается вперёд при каждом новом коммите: технически \texttt{main}~--- это файл \texttt{.git/refs/heads/main} с сорока символами внутри. \term{HEAD}~--- ссылка «где я сейчас», и обычно она указывает не на коммит, а на ветку; состояние, когда HEAD указывает прямо на коммит, называется \term{отсоединённым HEAD} (\eng{detached HEAD}), и сделанные в нём коммиты не принадлежат ни одной ветке. Запомните формулу: коммит~--- снимок, ветка~--- указатель на коммит, HEAD~--- указатель на текущую ветку.

\section{Три области и базовые команды}

Ключевая идея, которую нужно освоить в первую очередь: файл в Git живёт в трёх
областях, и переход между ними выполняется явными командами
(рис.~\ref{fig:12-2}).

\begin{figure}[htbp]\centering
\begin{tikzpicture}[font=\small]
\node[boxw=31mm] (wd) {Рабочий каталог\\[1pt]\scriptsize файлы на диске};
\node[boxw=33mm,right=25mm of wd] (idx) {Индекс\\[1pt]\scriptsize что войдёт в
  следующий коммит};
\node[boxw=31mm,right=25mm of idx] (rep) {Репозиторий\\[1pt]\scriptsize каталог
  \texttt{.git}};
\draw[flowarrow] (wd) -- node[above,font=\scriptsize] {\texttt{git add}} (idx);
\draw[flowarrow] (idx) -- node[above,font=\scriptsize] {\texttt{git commit}}
  (rep);
\draw[flowarrow] (rep.south) -- ++(0,-8mm) -| node[pos=0.25,below,
  font=\scriptsize] {\texttt{git restore}: вернуть файлы из истории} (wd.south);
\end{tikzpicture}
\caption{Три области, в которых живёт файл}\label{fig:12-2}
\end{figure}

Аналогия с супермаркетом объясняет схему за одну фразу: рабочий каталог~--- торговый зал, где вы берёте товары с полок и кладёте обратно; индекс~--- лента у кассы, куда выкладывают только то, что действительно собираются купить; коммит~--- пробитый чек. Именно индекс отличает Git от многих систем: он позволяет собрать аккуратный коммит из части изменений~--- вы поправили десять файлов, а в коммит кладёте три, относящихся к исправлению ошибки. Отсюда и четыре состояния файла: \term{неотслеживаемый} (\eng{untracked}), \term{изменённый} (\eng{modified}), \term{подготовленный} (\eng{staged}) и \term{зафиксированный} (\eng{committed}); команды, которыми файл проходит этот путь, собраны в листинге~\ref{lst:12-2}.

\begin{listing}[H]
\begin{minted}{bash}
git init -b main my-project   # новый репозиторий с веткой main
git clone <адрес>             # копия чужого репозитория

git status                    # что сейчас происходит
git add Book.java             # подготовить конкретный файл
git add -p                    # интерактивно, по кускам изменений

git commit -m "feat: искать книги по автору"
git commit --amend            # переделать последний коммит
\end{minted}
\caption{Создание репозитория и первые коммиты}
\label{lst:12-2}
\end{listing}

Команда \texttt{git clone} не просто «скачивает файлы»: она забирает всю базу объектов и привязывает локальную ветку к серверной. Команда \texttt{git add -p} показывает изменения кусками и спрашивает по каждому, включать ли его в коммит~--- лучший способ не утащить отладочный \texttt{System.out.println}, забытый в соседнем методе. А \texttt{git commit --amend} переписывает последний коммит, и у того меняется хеш: пока коммит не отправлен на сервер, это безопасно, а после отправки становится переписыванием опубликованной истории.

Сообщение коммита~--- письмо в будущее: себе через полгода и коллеге, который будет разбираться, почему строка выглядит именно так. Первая строка~--- до семидесяти символов, дальше, после пустой строки,~--- объяснение, \emph{почему} сделано именно так. Один коммит~--- одно логическое изменение, а префиксы \texttt{feat}, \texttt{fix}, \texttt{docs}, \texttt{test}, \texttt{refactor} читает и человек, и автоматика.

Историю читают командами из листинга~\ref{lst:12-3}. Ключ \texttt{--graph} рисует ветвления прямо в консоли, а \texttt{git blame} показывает хеш коммита для каждой строки файла: это не про поиск виноватых, а про поиск контекста.

\begin{listing}[H]
\begin{minted}{bash}
git log --oneline                    # по одной строке на коммит
git log --oneline --graph --all      # картинка ветвлений

git diff                 # рабочий каталог против индекса
git diff --staged        # индекс против последнего коммита
git show a3f1c9e         # содержимое конкретного коммита
git blame BookService.java   # кто написал каждую строку
\end{minted}
\caption{Просмотр истории}
\label{lst:12-3}
\end{listing}

В репозиторий попадает только исходный код и то, что нельзя сгенерировать: скомпилированные классы, JAR-файлы, настройки среды разработки и тем более пароли там не нужны. Список исключений задаёт файл \texttt{.gitignore} в корне проекта (листинг~\ref{lst:12-4}).

\begin{listing}[H]
\begin{minted}{text}
target/            # результаты сборки
out/
*.class
*.jar
.idea/             # настройки среды разработки
*.iml
*.log
application-local.properties   # локальные настройки
.env
\end{minted}
\caption{Типичный файл .gitignore для проекта на Java}
\label{lst:12-4}
\end{listing}

У этого файла есть нюанс, на котором спотыкаются все: правила действуют только на \emph{неотслеживаемые} файлы. Если файл уже попал в индекс, правило его оттуда не выкинет~--- нужна команда \texttt{git rm --cached application-local.properties}, убирающая файл из-под контроля версий, но оставляющая его на диске.

\begin{oshibka}
\textbf{Пароль в коммите.} Пароль, попавший в историю, считается скомпрометированным навсегда: даже если удалить его следующим коммитом, он останется в старом снимке и в клонах у всех, кто успел забрать изменения. Правильный порядок: сначала сменить пароль или ключ и только потом чистить историю.
\end{oshibka}

Отмена изменений~--- место, где теряются чаще всего. \texttt{Git restore} работает с файлами: с ключом \texttt{--staged} убирает файл из индекса, сохраняя правки, а без ключа выбрасывает правки необратимо. \texttt{Git reset} двигает указатель ветки и, в зависимости от режима, трогает индекс и рабочий каталог (табл.~\ref{tab:12-2}). \texttt{Git revert} ничего не удаляет, а создаёт новый коммит с обратными изменениями.

\begin{table}[htbp]
\caption{Режимы команды \texttt{git reset}}
\label{tab:12-2}
\centering
\begin{tabular}{L{2.6cm}L{2.2cm}L{2.6cm}L{7cm}}
\toprule
Режим & Индекс & Рабочий каталог & Когда применять \\
\midrule
\texttt{--soft} & не тронут & не тронут & переделать последний коммит: изменения остаются подготовленными \\
\texttt{--mixed} & сброшен & не тронут & разложить последний коммит на несколько \\
\texttt{--hard} & сброшен & перезаписан & выбросить последний коммит вместе с изменениями \\
\bottomrule
\end{tabular}
\end{table}

Режим \texttt{--hard}~--- единственная по-настоящему опасная команда из списка: незакоммиченные правки она уничтожает без вопросов, и \texttt{reflog} их не вернёт, потому что помнит коммиты, а не правки. Отсюда правило: историю, которую видели другие, отменяют через \texttt{git revert}, а \texttt{git reset} применяют только к своим локальным коммитам.

\section{Ветвление, слияние и конфликты}

Ветка в Git~--- файл с хешем коммита, поэтому её создание не копирует ни одного файла проекта и выполняется мгновенно. Ветка~--- закладка в книге: сколько бы закладок вы ни положили, книга не станет толще. Смысл ветки~--- изоляция: пока вы делаете поиск в ветке \texttt{feature/search}, ветка \texttt{main} остаётся пригодной для выпуска. Команды работы с ветками собраны в листинге~\ref{lst:12-5}.

\begin{listing}[H]
\begin{minted}{bash}
git branch -a                  # список веток, включая серверные
git switch -c feature/search   # создать и переключиться
git switch main                # переключиться на существующую
git branch -d feature/search   # удалить слитую ветку
git branch -D feature/draft    # удалить принудительно
\end{minted}
\caption{Работа с ветками}
\label{lst:12-5}
\end{listing}

Исторически для переключения использовалась команда \texttt{git checkout}, которая делает слишком много разного сразу; с версии~2.23 её разделили: \texttt{switch} переключает ветки, \texttt{restore} восстанавливает файлы. Имена веток задают соглашением: \texttt{feature/book-search}, \texttt{bugfix/npe-in-search}, \texttt{release/1.4.0}.

Простейший случай слияния~--- \term{быстрая перемотка} (\eng{fast-forward}): вы создали ветку от \texttt{main}, сделали два коммита, а \texttt{main} за это время не менялась. Прямой путь вперёд есть, поэтому Git просто перемещает указатель \texttt{main} на вершину ветки, не создавая нового коммита. История остаётся линейной, но факт существования отдельной ветки из неё исчезает; если он важен, слияние делают ключом \texttt{--no-ff}, и тогда каждая задача видна в истории отдельным «пузырём».

Реальная ситуация сложнее: пока вы работали в своей ветке, в \texttt{main} тоже появились коммиты (верхний ряд рис.~\ref{fig:12-3}). Прямого пути вперёд нет, и Git выполняет \term{трёхстороннее слияние} (\eng{three-way merge}): берёт вершину \texttt{main}, вершину \texttt{feature} и их общего предка~--- базу слияния~--- и сравнивает их. Если файл изменился относительно базы только в одной ветке, берётся изменённая версия; если в обеих, но в разных местах~--- изменения объединяются; если изменились одни и те же строки~--- возникает конфликт, а результатом становится \term{коммит слияния} с двумя родителями (нижний ряд рисунка).

\begin{figure}[htbp]\centering
\begin{tikzpicture}[font=\footnotesize,
  cnode/.style={draw,circle,inner sep=0pt,minimum size=8mm,font=\scriptsize}]
\node[font=\scriptsize\itshape] (l1) at (0,-0.7) {до слияния};
\node[cnode] (a1) at (2.2,0) {C1};
\node[cnode] (a2) at (3.5,0) {C2};
\node[cnode] (a5) at (4.8,0) {C5};
\node[cnode] (a6) at (6.1,0) {C6};
\node[cnode] (a3) at (4.8,-1.4) {C3};
\node[cnode] (a4) at (6.1,-1.4) {C4};
\draw[-Latex] (a1) -- (a2);
\draw[-Latex] (a2) -- (a5);
\draw[-Latex] (a5) -- (a6);
\draw[-Latex] (a2) -- (a3);
\draw[-Latex] (a3) -- (a4);
\node[font=\scriptsize] (n1) at (7.4,0) {\texttt{main}};
\node[font=\scriptsize] (n2) at (7.5,-1.4) {\texttt{feature}};
\node[font=\scriptsize\itshape] (l2) at (0,-3.6) {после слияния};
\node[cnode] (b1) at (2.2,-2.9) {C1};
\node[cnode] (b2) at (3.5,-2.9) {C2};
\node[cnode] (b5) at (4.8,-2.9) {C5};
\node[cnode] (b6) at (6.1,-2.9) {C6};
\node[cnode] (b3) at (4.8,-4.3) {C3};
\node[cnode] (b4) at (6.1,-4.3) {C4};
\node[cnode,very thick] (b7) at (7.4,-2.9) {C7};
\draw[-Latex] (b1) -- (b2);
\draw[-Latex] (b2) -- (b5);
\draw[-Latex] (b5) -- (b6);
\draw[-Latex] (b2) -- (b3);
\draw[-Latex] (b3) -- (b4);
\draw[-Latex] (b6) -- (b7);
\draw[-Latex] (b4) -- (b7);
\node[font=\scriptsize] (n3) at (8.7,-2.9) {\texttt{main}};
\end{tikzpicture}
\caption{Трёхстороннее слияние и коммит C7 с двумя родителями}\label{fig:12-3}
\end{figure}

\term{Конфликт} возникает, когда две ветки изменили одни и те же строки одного файла: Git не может решить за вас, чей вариант правильный. Это не поломка, а нормальная ситуация~--- два повара дописали в одну и ту же строку рецепта «варить 20~минут» и «варить 40~минут», и автоматика не знает, какой суп вы хотите. Git оставляет файл с \term{маркерами конфликта} (листинг~\ref{lst:12-6}).

\begin{listing}[H]
\begin{minted}{java}
public List<Book> findBooks(String query) {
<<<<<<< HEAD
    // Вариант из ветки main: поиск только по названию
    return repository.findByTitle(query);
=======
    // Вариант из feature/search: по названию и автору
    return repository.findByTitleOrAuthor(query, query);
>>>>>>> feature/search
}
\end{minted}
\caption{Маркеры конфликта в файле}
\label{lst:12-6}
\end{listing}

Между первым маркером и \texttt{=======} лежит версия текущей ветки, той, в которую вы сливаете; ниже~--- версия ветки, которую сливаете. Разрешают конфликт в четыре шага: \texttt{git status} перечисляет конфликтующие файлы; вы пишете правильный вариант~--- чаще всего третий, объединяющий оба намерения; убираете все маркеры (забытый маркер~--- ошибка компиляции); собираете проект и прогоняете тесты, потому что «склеенный» код часто собирается, но делает не то (листинг~\ref{lst:12-7}).

\begin{listing}[H]
\begin{minted}{bash}
git status                 # список конфликтующих файлов
git add BookService.java   # пометить конфликт разрешённым
git commit --no-edit       # принять сообщение, готовое у Git

git merge --abort          # отказаться от слияния целиком
git checkout --ours   BookService.java   # версия своей ветки
git checkout --theirs BookService.java   # версия чужой ветки
\end{minted}
\caption{Команды при разрешении конфликта}
\label{lst:12-7}
\end{listing}

Ключи \texttt{--ours} и \texttt{--theirs} берут файл целиком, поэтому правки второй стороны теряются. Чтобы конфликтовать реже, делайте маленькие задачи и короткоживущие ветки, регулярно подтягивайте изменения из основной ветки и договоритесь о правилах форматирования: половина конфликтов~--- это войны табуляций и пробелов.

\section{Удалённые репозитории, ревью и модели ветвления}

\term{Удалённый репозиторий} (\eng{remote})~--- другая копия того же репозитория, обычно на сервере; Git хранит для неё короткое имя, по умолчанию \texttt{origin}. У каждого разработчика дома полная мастерская, а \texttt{origin}~--- общий склад, куда свозят готовые детали. Команды обмена приведены в листинге~\ref{lst:12-8}.

\begin{listing}[H]
\begin{minted}{bash}
git remote add origin <адрес сервера>

git fetch origin                    # скачать, не трогая файлы
git log --oneline HEAD..origin/main # что нового на сервере

git pull origin main                # fetch + merge одной командой
git push origin main                # отправить коммиты
git push -u origin feature/search   # первая отправка со связью
\end{minted}
\caption{Обмен с удалённым репозиторием}
\label{lst:12-8}
\end{listing}

Разница между \texttt{fetch} и \texttt{pull} принципиальна. Первая обновляет ваши знания о сервере: появляются ветки вида \texttt{origin/main}, вы смотрите, что изменилось, и решаете, что делать; рабочий каталог не меняется. Вторая~--- это \texttt{fetch} плюс слияние одной командой: удобно, но конфликт свалится на вас в неподходящий момент. \term{Отслеживаемая ветка} (\eng{tracking branch})~--- локальная ветка, связанная с серверной; связь устанавливает ключ \texttt{-u}, и благодаря ей отправка и получение работают без аргументов.

\begin{oshibka}
\textbf{Отклонённая отправка.} Сообщение \texttt{rejected~--- non-fast-forward} означает, что на сервере есть коммиты, которых нет у вас. Правильное решение~--- забрать их и отправить снова; неправильное~--- \texttt{git push --force}, затирающий чужие коммиты. Если принудительная отправка действительно нужна, берите \texttt{git push --force-with-lease}: она откажется затирать ветку, если на сервере появилось что-то, чего вы не видели.
\end{oshibka}

Прямая запись в основную ветку в командах почти всегда запрещена: вместо неё используется механизм предложения изменений~--- \term{pull request} на одних платформах и \term{merge request} на других. Вы не вклеиваете свою главу в общую монографию сами, а приносите её редактору. Жизненный цикл задачи показан в листинге~\ref{lst:12-9}; хорошее описание говорит, что сделано, зачем и как проверить результат.

\begin{listing}[H]
\begin{minted}{bash}
git switch main && git pull origin main   # 1. обновиться
git switch -c feature/book-search         # 2. ветка под задачу
git commit -am "feat: добавить поиск книг по автору"
git push -u origin feature/book-search    # 3. отправить ветку
# 4. открыть pull request в веб-интерфейсе платформы
# 5. дождаться зелёного конвейера и одобрения ревьюеров
git switch main && git pull origin main   # 6. после слияния
git branch -d feature/book-search         #    удалить ветку
\end{minted}
\caption{Жизненный цикл одной задачи}
\label{lst:12-9}
\end{listing}

\term{Код-ревью} (\eng{code review})~--- обязательный просмотр изменений другим разработчиком до попадания кода в основную ветку: он находит ошибки до продуктива, распространяет знание о системе и поддерживает единый стиль. Комментируйте код, а не человека («здесь возможен \texttt{NullPointerException}, если список пуст» вместо «ты опять не подумал»), и делайте маленькие предложения изменений: правка на двести строк получит осмысленное ревью, правка на три тысячи~--- комментарий «выглядит нормально». Механические проверки выполняет конвейер: ревьюер слишком дорогой линтер. Если прав на запись нет, делают \term{форк} (\eng{fork})~--- копию чужого репозитория, добавляя оригинал вторым удалённым под именем \texttt{upstream}.

Команда из десяти человек не может создавать ветки как придётся~--- нужна договорённость о том, какие ветки существуют и куда сливаются. Сложились три модели. \term{Git Flow}~--- постоянные \texttt{main} и \texttt{develop} плюс временные \texttt{feature/*}, \texttt{release/*} и \texttt{hotfix/*}: строгий порядок для продукта с несколькими поддерживаемыми версиями, но тяжеловесно. \term{GitHub Flow}~--- одна постоянная работоспособная ветка, ветка на задачу, ревью, слияние: лучший выбор для веб-сервиса и учебных проектов. \term{Trunk-based}~--- все работают в основной ветке, незавершённое прячут за флагами функций: конфликтов почти нет, но без сильных тестов сломанная основная ветка блокирует команду. Принцип общий: чем короче живёт ветка, тем меньше боли при слиянии.

\section{Перебазирование и повседневные инструменты}

\term{Перебазирование} (\eng{rebase}) переносит коммиты вашей ветки так, будто вы начали работу не от старой точки, а от актуальной вершины основной ветки. Возьмите верхний ряд рис.~\ref{fig:12-3} за исходное состояние: после команды \texttt{git rebase main}, выполненной в ветке \texttt{feature}, коммиты C3 и C4 переносятся за C6, и история вытягивается в одну линию~--- C1, C2, C5, C6, C3', C4'. Штрихи не случайны: C3' и C4'~--- новые коммиты. Изменения те же, но родитель другой, значит и хеш другой; старые коммиты остаются в базе объектов, но ни одна ветка на них уже не указывает.

Слияние~--- вклеить в тетрадь дополнительный лист с пометкой «здесь сошлись два черновика»; перебазирование~--- переписать свои страницы набело сразу после последней страницы соседа: аккуратнее, но переписанные страницы уже не те, что были. Обе операции решают одну задачу, но платят по-разному. Слияние сохраняет реальный ход событий и хеши коммитов, добавляя коммит слияния; оно безопасно всегда, и конфликт разрешается один раз. Перебазирование даёт линейную историю без лишнего коммита, но создаёт новые коммиты, опасно для опубликованных веток и может потребовать разрешать конфликт на каждом переносимом коммите.

Отсюда золотое правило: \emph{никогда не перебазируйте ветку, которую видели другие}~--- иначе коллега получит конфликты в коммитах, которых не писал. Личная ветка~--- перебазировать можно; общая~--- только слияние.

Реальная история личной ветки редко выглядит красиво: «доделал», «фикс», «ещё
раз фикс». Перед тем как показывать её людям, историю приводят в порядок
командой \texttt{git rebase -i HEAD\textasciitilde 4}: откроется редактор со
списком коммитов от старых к новым, где первое слово каждой строки~--- команда
(листинг~\ref{lst:12-10}).

\begin{listing}[H]
\begin{minted}{text}
pick   8f2c1a4 добавил поиск по автору
squash 3d7f4b1 ещё раз фикс
squash 5e0c2d8 фикс
reword 71ba9f3 доделал
\end{minted}
\caption{План интерактивного перебазирования}
\label{lst:12-10}
\end{listing}

Команд шесть: \texttt{pick} оставляет коммит как есть, \texttt{reword} меняет только сообщение, \texttt{edit} останавливает перебазирование для правки коммита, \texttt{squash} и \texttt{fixup} объединяют коммит с предыдущим (второй выбрасывает своё сообщение), \texttt{drop} удаляет коммит. В нашем плане две строки \texttt{squash} вливаются в стоящий выше \texttt{pick}, а \texttt{reword} оставляет последний коммит отдельным: из четырёх коммитов останется два. Полезен и \texttt{git pull --rebase}: он переносит поверх чужих изменений ваши локальные коммиты вместо того, чтобы плодить коммиты «Merge branch...».

Ежедневная работа опирается ещё на три команды (листинг~\ref{lst:12-11}). \texttt{Git stash} прячет незаконченные правки, когда посреди задачи просят срочно посмотреть другую ветку: это верхний ящик стола, куда смахивают бумаги и достают их потом ровно в том же виде. \term{Тег}~--- неподвижная метка на коммите; для релизов берут аннотированный тег, а имена следуют \term{семантическому версионированию} (\eng{SemVer}): мажорная часть меняется при несовместимых изменениях, минорная при добавлении функциональности, патч при исправлениях. \texttt{Git cherry-pick} переносит отдельный коммит в другую ветку, но коммит получает новый хеш, поэтому основным способом переноса кода быть не должен.

\begin{listing}[H]
\begin{minted}{bash}
git stash push -m "недоделанный поиск"   # спрятать правки
git stash pop                # достать и убрать из списка

git tag -a v1.0.0 -m "Первый релиз библиотеки"
git push origin --tags       # теги не уходят автоматически

git switch release/1.4
git cherry-pick a3f1c9e      # перенести один коммит
\end{minted}
\caption{Отложенная работа, теги и перенос коммита}
\label{lst:12-11}
\end{listing}

Самая успокаивающая команда Git~--- \texttt{git reflog}: журнал всех перемещений HEAD и веток в локальном репозитории, корзина операционной системы, из которой удалённое можно достать, пока её не очистили. Команда \texttt{git bisect} решает задачу «ошибка появилась где-то за последние триста коммитов»: двоичный поиск укладывается в восемь-девять шагов, как в игре «угадай число от 1 до 1000 за десять вопросов». В автоматическом режиме (листинг~\ref{lst:12-12}) критерием служит код возврата: ноль~--- «хорошо», иначе~--- «плохо», а код~125 означает «коммит проверить невозможно», и такой коммит пропускается.

\begin{listing}[H]
\begin{minted}{bash}
git reflog                        # журнал перемещений HEAD
git reset --hard HEAD@{1}         # вернуть ветку до сброса

git bisect start
git bisect bad                    # текущее состояние сломано
git bisect good v1.0.0            # здесь всё работало
git bisect run mvn -q test        # автоматический поиск
git bisect reset                  # вернуться в исходное
\end{minted}
\caption{Восстановление истории и поиск виновного коммита}
\label{lst:12-12}
\end{listing}

Наконец, \term{Git LFS} (\eng{Large File Storage}) подменяет большие двоичные файлы текстовыми указателями и держит содержимое на отдельном сервере: сам Git хранит каждую версию такого файла целиком и не умеет их сливать.

\section{Непрерывная интеграция и доставка}

Классическая сцена: у разработчика всё собирается, а у коллеги проект не компилируется, потому что не хватает файла, не попавшего в коммит; или тесты, которые никто не запускал две недели, падают все разом. Фраза «но у меня на машине работает»~--- симптом отсутствия непрерывной интеграции. Лекарство~--- заводской конвейер: деталь на ленте автоматически проходит проверку, и брак отсеивается сразу.

Три близких понятия различают так. \term{Непрерывная интеграция} (\eng{Continuous Integration, CI}) после каждого изменения автоматически собирает проект и прогоняет тесты в чистом окружении. \term{Непрерывная доставка} (\eng{Continuous Delivery, CD}) дополнительно готовит проверенный артефакт к выпуску, но решение о выкладке принимает человек. \term{Непрерывное развёртывание} (\eng{Continuous Deployment}) выкладывает каждое прошедшее проверки изменение само: разница с доставкой ровно в одном~--- есть ли кнопка, которую нажимает человек. Выгода видна сразу: ошибку находят через пять минут после коммита, сборка идёт в чистом контейнере, а основную ветку защищает запрет сливать изменения, пока конвейер красный.

Устроен конвейер из нескольких понятий. \term{Конвейер} (\eng{pipeline})~--- весь процесс, запускаемый одним событием: отправкой коммитов, предложением изменений, появлением тега. \term{Стадия} (\eng{stage})~--- группа заданий, например \texttt{build}, \texttt{test}, \texttt{deploy}; \term{задание} (\eng{job})~--- набор команд; \term{раннер} (\eng{runner})~--- агент, исполняющий задание в контейнере; \term{артефакты}~--- файлы, сохраняемые после задания; \term{кеш}~--- данные, переиспользуемые между запусками. Стадии выполняются последовательно, задания внутри стадии~--- параллельно; если задание упало, следующие стадии не запускаются. Как это выглядит для проекта на Maven, показано в листинге~\ref{lst:12-13}.

\begin{listing}[H]
\begin{minted}{text}
default:
  image: maven:3.9-eclipse-temurin-21   # Maven и JDK 21

variables:
  MAVEN_OPTS: "-Dmaven.repo.local=$CI_PROJECT_DIR/.m2/repository"

cache:                    # не качать зависимости заново
  key: maven-repository
  paths:
    - .m2/repository

stages: [build, test, deploy]

compile:
  stage: build
  script:
    - mvn --batch-mode compile

unit-test:
  stage: test
  script:
    - mvn --batch-mode test
  artifacts:
    when: always          # отчёты нужны и при упавших тестах
    reports:
      junit:
        - target/surefire-reports/TEST-*.xml

deploy-production:
  stage: deploy
  script:
    - echo "Выкладываем релиз $CI_COMMIT_TAG"
  rules:                  # только по тегу и только по кнопке
    - if: $CI_COMMIT_TAG =~ /^v\d+\.\d+\.\d+$/
      when: manual
\end{minted}
\caption{Конвейер сборки, тестирования и выкладки}
\label{lst:12-13}
\end{listing}

Раздел \texttt{image} задаёт образ контейнера: он уже содержит Maven и JDK. Переменная \texttt{MAVEN\_OPTS} перенаправляет локальный репозиторий Maven внутрь каталога проекта, потому что домашний каталог между заданиями не сохраняется и потому не кешируется; без кеша каждый запуск качал бы десятки мегабайт заново. Артефакт вида \texttt{artifacts:reports:junit} платформа разбирает сама и показывает результаты тестов в предложении изменений, а \texttt{artifacts:paths} просто сохраняет файлы. Ключ \texttt{when: manual} превращает задание в кнопку; переменные с префиксом \texttt{CI\_} подставляет платформа.

\begin{vrezka}
\textbf{Секреты в конвейере.} Пароли, токены и ключи никогда не пишут в файле конвейера: он лежит в репозитории и доступен каждому, у кого есть клон. Их задают в настройках проекта как защищённые переменные CI/CD, откуда они попадают в задание как обычные переменные окружения.
\end{vrezka}

Вторая распространённая система описывает конвейеры файлами в каталоге \texttt{.github/workflows}: синтаксис другой, идея та же~--- задание запускается событием \texttt{on: push} или \texttt{on: pull\_request}, состоит из шагов, и шаг может переиспользовать готовое действие (\texttt{actions/checkout}, \texttt{actions/setup-java}) вместо собственных команд. Освоив одну систему, на другую переходят за день.

Три правила делают конвейер полезным, а не декоративным: он должен быть быстрым (ориентир~--- до десяти минут, иначе его перестают ждать); красная основная ветка чинится первой; тесты не должны зависеть от порядка запуска и внешних сервисов, потому что «плавающие» тесты приучают перезапускать задание вместо чтения ошибки.

\section{Гибкие процессы: Agile, Scrum и Kanban}

\term{Каскадная модель} (\eng{Waterfall}) предполагает строгую последовательность этапов~--- требования, проектирование, реализация, тестирование, внедрение,~--- где каждый начинается только после предыдущего. Она работает там, где требования известны заранее, но в разработке программ регулярно давала сбой: требования устаревали за время реализации, заказчик видел результат только в конце, а ошибка в требованиях, найденная при тестировании, стоила дороже всего. Это как заказать ужин по телефону и увидеть блюдо через три часа, когда менять что-либо поздно.

В 2001~году семнадцать разработчиков сформулировали \term{Манифест гибкой разработки программного обеспечения} (\eng{Agile Manifesto}) из четырёх ценностей: люди и взаимодействие важнее процессов и инструментов; работающий продукт важнее исчерпывающей документации; сотрудничество с заказчиком важнее согласования условий контракта; готовность к изменениям важнее следования первоначальному плану. Ключевая оговорка часто теряется при пересказе: не отрицая важности того, что справа, авторы всё-таки больше ценят то, что слева~--- подход не отменяет документацию, планы и договоры, а расставляет приоритеты. Из двенадцати принципов стоит выделить два: работающий продукт назван основным показателем прогресса (отсюда важность конвейера), а внимание к техническому совершенству~--- условием гибкости, то есть «быстро и как получится» подходу противоречит.

\term{Scrum}~--- лёгкий фреймворк для решения сложных задач с изменяющимися требованиями: он описывает не то, как писать код, а то, как команда организует работу. Название пришло из регби: схватка, в которой команда двигает мяч вперёд общим усилием~--- задача не «делится по людям», а берётся командой. В основе~--- прозрачность, инспекция и адаптация, а структура укладывается в формулу «три роли, три артефакта, пять событий».

Ролей три. \term{Владелец продукта} (\eng{Product Owner}) отвечает за ценность продукта, содержание и порядок бэклога и связь с заказчиком, но не раздаёт задачи. \term{Скрам-мастер} (\eng{Scrum Master}) отвечает за работоспособность процесса и устранение препятствий, но не является начальником. \term{Команда разработчиков} создаёт инкремент, оценивает задачи и отвечает за качество. Владелец продукта решает «что», команда~--- «как» и «сколько успеем».

Артефактов тоже три. \term{Бэклог продукта}~--- упорядоченный по ценности список всего, что может понадобиться продукту; за него отвечает владелец продукта. \term{Бэклог спринта}~--- задачи, выбранные на текущий спринт, вместе с его целью. \term{Инкремент}~--- работоспособный результат спринта. Элементы бэклога записывают в виде \term{пользовательских историй} (\eng{user stories}), образец~--- в листинге~\ref{lst:12-14}: формулировка от лица пользователя возвращает команду к вопросу «зачем это нужно человеку».

\begin{listing}[H]
\begin{minted}{text}
Как читатель библиотеки,
я хочу искать книги по фамилии автора,
чтобы находить издание, не зная точного названия.

Критерии приёмки:
- поиск не учитывает регистр;
- при пустом результате показывается понятное сообщение;
- ответ приходит быстрее 300 мс на базе в 10 000 книг.
\end{minted}
\caption{Пользовательская история с критериями приёмки}
\label{lst:12-14}
\end{listing}

Событий пять. \term{Спринт}~--- отрезок в одну-четыре недели, за который создаётся инкремент; длительность его постоянна, а цель не меняется по ходу. \term{Планирование} (до четырёх часов) формулирует цель и отбирает элементы бэклога, \term{ежедневный скрам} (пятнадцать минут) синхронизирует команду, \term{обзор спринта} (до двух часов) показывает инкремент, а \term{ретроспектива} (до полутора часов) разбирает процесс и заканчивается конкретными действиями.

Обзор и ретроспективу часто путают: обзор~--- про продукт, ретроспектива~--- про процесс, и ретроспектива без конкретных действий с ответственным~--- потерянное время. \term{Определение готовности} (\eng{Definition of Done})~--- общий список условий, при которых работа считается завершённой, чтобы слово «готово» означало одно и то же для всех: код написан и покрыт тестами, прошёл ревью, конвейер зелёный, изменения слиты в основную ветку.

Оценивают работу не в часах. \term{Стори-поинты}~--- безразмерная оценка объёма: не «сколько часов», а «насколько эта задача больше или меньше вон той»; обычно берут ряд Фибоначчи 1, 2, 3, 5, 8, 13. Причина проста: люди систематически ошибаются в абсолютных оценках времени, но неплохо сравнивают задачи между собой. \term{Скорость} (\eng{velocity})~--- среднее число поинтов за спринт; она нужна только для планирования следующего спринта: требовать её роста~--- прямой путь к завышению оценок. Ход спринта показывает \term{диаграмма сгорания} (\eng{burndown chart})~--- график оставшейся работы по дням: если фактическая линия долго идёт горизонтально, работа не завершается, хотя «почти готова».

\term{Kanban}~--- альтернативный подход, пришедший из производственной системы Toyota: спринтов, фиксированных ролей и обязательных событий здесь нет, есть поток задач и правила управления им~--- визуализировать работу на доске со столбцами по стадиям процесса, ограничивать незавершённую работу \term{WIP-лимитом} (максимальным числом карточек в столбце), измерять время прохождения задачи и делать правила перехода явными. Смысл WIP-лимита понятен на кухне с четырьмя конфорками: восемь кастрюль поставить можно, но еда не приготовится быстрее, зато что-нибудь подгорит. Отличия от Scrum: вместо спринтов~--- непрерывный поток, обязательных ролей нет, а ключевая метрика~--- не скорость команды, а время прохождения задачи.

Задачи ведут в трекерах, и важнее выбора инструмента~--- связь трекера с репозиторием: строка вида \texttt{Closes \#42} в сообщении коммита заставит платформу после слияния закрыть задачу номер~42. Это самая дешёвая проектная документация: она пишется сама. Так и связаны три темы главы: Git даёт технику совместной работы, CI/CD~--- автоматизацию проверок, Scrum~--- ритм и обратную связь.

\praktikum{}

Понадобятся Git версии~2.23 или новее (именно в ней появились \texttt{switch} и \texttt{restore}), JDK~21 и Maven. Работайте в командной строке: пока вы не увидите, что происходит с указателями вручную, графический интерфейс будет казаться магией. Выполните настройку из листинга~\ref{lst:12-1} и обязательно задайте редактор (\texttt{git config --global core.editor "notepad"} на Windows, \texttt{"nano"} на Linux и macOS): выйти из незнакомого редактора~--- отдельный навык (в nano это Ctrl+O, Enter, Ctrl+X; в vim~--- Esc и команда \texttt{:wq}).

\subsection*{Задание 1. Репозиторий, история и файл .gitignore}

Создайте каталог для экспериментов и выполните \texttt{git init -b main library}. Внутри создайте файл \texttt{README.md} со строкой «Библиотека» и проследите три состояния файла: \texttt{git status} до и после \texttt{git add}, затем коммит \texttt{docs: добавить README}. Создайте каталог \texttt{src/main/java/ru/fa/library}, положите туда классы из листинга~\ref{lst:12-15} в кодировке UTF-8, проверьте сборку командой \texttt{javac -d out -sourcepath src/main/java} и запуск командой \texttt{java -cp out ru.fa.library.LibraryApp}, а затем сделайте \emph{три отдельных} коммита~--- по одному на класс.

\begin{listing}[H]
\begin{minted}{java}
package ru.fa.library;   // файл Book.java

public record Book(String title, String author, int year) {
}

// ---------- файл BookService.java ----------
package ru.fa.library;

import java.util.ArrayList;
import java.util.List;

public class BookService {

    private final List<Book> books = new ArrayList<>();

    public void add(Book book) {
        books.add(book);
    }

    public List<Book> findAll() {
        return List.copyOf(books);
    }
}

// ---------- файл LibraryApp.java ----------
package ru.fa.library;

public class LibraryApp {

    public static void main(String[] args) {
        BookService service = new BookService();
        service.add(new Book("Война и мир", "Толстой", 1869));
        service.add(new Book("Мастер и Маргарита", "Булгаков", 1967));
        service.findAll().forEach(book ->
                System.out.println(book.title()));
    }
}
\end{minted}
\caption{Классы учебного проекта}
\label{lst:12-15}
\end{listing}

Добавьте в \texttt{BookService} метод \texttt{count()}, возвращающий \texttt{books.size()}, и, \emph{не коммитя сразу}, сравните вывод команд \texttt{git diff} и \texttt{git diff --staged} до и после \texttt{git add}; закоммитьте изменение и разберите вывод команд из листинга~\ref{lst:12-3}. Затем создайте \texttt{.gitignore} по образцу листинга~\ref{lst:12-4} и убедитесь, что каталог \texttt{out} исчез из вывода \texttt{git status}. Наконец, смоделируйте типичную ошибку: закоммитьте файл \texttt{application-local.properties} со строкой \texttt{db.password=SuperSecret123}, добавьте его имя в \texttt{.gitignore}, убедитесь, что правило не сработало, и исправьте это командой \texttt{git rm --cached application-local.properties}.

\emph{Ответьте письменно:} что означает состояние \texttt{untracked}; почему три коммита лучше одного «добавил всё»; почему правило \texttt{.gitignore} не подействовало на уже отслеживаемый файл; что и в каком порядке делают, когда пароль попал в историю.

\subsection*{Задание 2. Ветки, слияние и конфликт}

Создайте ветку \texttt{feature/find-by-title}, добавьте в \texttt{BookService} метод \texttt{findBooks}, отбирающий книги, чьё название содержит подстроку запроса, и закоммитьте его; вторым коммитом добавьте вызов поиска в \texttt{LibraryApp}. Вернитесь на \texttt{main} и выполните слияние: убедитесь, что в выводе есть слово \texttt{Fast-forward} и нового коммита не появилось. Повторите упражнение с веткой \texttt{feature/oldest-book} и методом \texttt{oldestYear()}, который возвращает минимальный год издания, а для пустой библиотеки~--- ноль; слейте её командой \texttt{git merge --no-ff} и сравните графы двух слияний командой \texttt{git log --oneline --graph -8}.

\begin{sloppypar}
Теперь создайте конфликт намеренно. В ветке \texttt{feature/search-by-author} перепишите тело \texttt{findBooks} так, чтобы поиск шёл по фамилии автора; вернитесь на \texttt{main} и перепишите \emph{тот же} метод иначе~--- поиск по названию без учёта регистра. Закоммитьте оба варианта и выполните слияние: Git сообщит о конфликте. Осмотритесь, не разрешая его, и отмените слияние командой \texttt{git merge --abort}. Повторите слияние и разрешите конфликт по-настоящему: удалите все три маркера и напишите третий, правильный вариант~--- поиск и по названию, и по фамилии автора, без учёта регистра. Убедитесь, что маркеров не осталось, соберите и запустите проект и лишь затем завершите слияние командами \texttt{git add} и \texttt{git commit --no-edit}.
\end{sloppypar}

\emph{Ответьте письменно:} что означает \texttt{Fast-forward} и куда переместился указатель; чем граф после \texttt{--no-ff} отличается от предыдущего; почему Git не смог слить изменения автоматически, хотя правки в других файлах объединил без вопросов; почему после разрешения конфликта обязательно собирать проект.

\subsection*{Задание 3. Слияние против перебазирования и восстановление истории}

Чтобы сравнить два способа честно, нужны две одинаковые ветки. Создайте рядом репозиторий \texttt{rebase-demo}, сделайте в нём два коммита в основной ветке, затем ветку \texttt{feature/report} ещё с двумя коммитами и скопируйте её командой \texttt{git branch feature/report-copy}. Вернитесь в основную ветку, добавьте туда ещё два коммита и запомните состояние командой \texttt{git branch stable}. Одну ветку слейте, другую перебазируйте на \texttt{stable} и сравните результат командой \texttt{git log --oneline --graph --all}.

Вернитесь в репозиторий \texttt{library}, создайте ветку \texttt{feature/statistics} и добавьте в \texttt{BookService} метод \texttt{averageYear()}, возвращающий средний год издания (для пустой библиотеки~--- ноль). Сделайте четыре коммита с намеренно плохими сообщениями \texttt{wip}, \texttt{фикс}, \texttt{ещё раз фикс}, \texttt{готово} и приведите историю в порядок командой \texttt{git rebase -i HEAD\textasciitilde 4} по образцу листинга~\ref{lst:12-10}, объединив их в один коммит \texttt{feat: добавить средний год издания}. Слейте ветку в основную.

Наконец, сломайте историю и почините её. Выпишите вывод \texttt{git log --oneline -5}, выполните \texttt{git reset --hard HEAD\textasciitilde 3} и убедитесь, что три коммита исчезли. Найдите прежнюю позицию в выводе \texttt{git reflog} и верните ветку командой \texttt{git reset --hard HEAD@\{1\}}; затем удалите ветку \texttt{feature/statistics} принудительно и восстановите её по хешу из reflog.

\emph{Ответьте письменно:} сколько коммитов оказалось в каждой из двух веток и откуда взялась разница; совпадают ли хеши коммитов с одинаковым содержимым и почему; какая история читается удобнее, а какая честнее отражает ход событий; что хранит reflog и когда он не поможет.

\subsection*{Задание 4. Удалённый репозиторий и конвейер}

Аккаунт на внешней платформе для первой части задания не нужен: сервером будет обычный каталог. Выполните \texttt{git init --bare ../library-server.git}, привяжите его как \texttt{origin} и отправьте ветку и теги. Изобразите коллегу: склонируйте репозиторий во второй каталог, допишите строку в \texttt{README.md} и отправьте изменение оттуда. Вернитесь в свой репозиторий, создайте файл \texttt{CONTRIBUTING.md}, сделайте коммит и попробуйте отправить его, не забирая чужие изменения~--- отправка будет отклонена. Разберитесь по-правильному: \texttt{git fetch origin}, \texttt{git pull --rebase origin main}, повторная отправка.

Переведите проект на Maven: создайте в корне файл \texttt{pom.xml} с идентификаторами \texttt{ru.fa} и \texttt{library}, версией~1.0.0, свойством \texttt{maven.compiler.release} равным~21, зависимостью \texttt{org.junit.jupiter:junit-jupiter} версии 5.11.4 в области видимости \texttt{test} и плагином \texttt{maven-surefire-plugin} версии 3.5.2. Добавьте каталог \texttt{src/test/java/ru/fa/library} и тест из листинга~\ref{lst:12-16}, затем соберите проект командами \texttt{mvn -q clean test} и \texttt{mvn package}.

\begin{listing}[H]
\begin{minted}{java}
package ru.fa.library;

import org.junit.jupiter.api.BeforeEach;
import org.junit.jupiter.api.Test;

import static org.junit.jupiter.api.Assertions.assertEquals;
import static org.junit.jupiter.api.Assertions.assertTrue;

class BookServiceTest {

    private final BookService service = new BookService();

    @BeforeEach
    void fillLibrary() {
        service.add(new Book("Война и мир", "Толстой", 1869));
        service.add(new Book("Мастер и Маргарита", "Булгаков", 1967));
    }

    @Test
    void findsBookByAuthor() {
        assertEquals(1, service.findBooks("булгаков").size());
    }

    @Test
    void handlesEmptyLibrary() {
        BookService empty = new BookService();
        assertEquals(0.0, empty.averageYear());
        assertTrue(empty.findBooks("Пушкин").isEmpty());
    }
}
\end{minted}
\caption{Тесты службы книг}
\label{lst:12-16}
\end{listing}

Добавьте в корень репозитория файл конвейера из листинга~\ref{lst:12-13} (на платформе с каталогом \texttt{.github/workflows} опишите те же стадии её синтаксисом), закоммитьте и отправьте его. Если у вас есть аккаунт на внешней платформе, создайте там \emph{пустой} репозиторий, привяжите его вторым удалённым и дождитесь завершения конвейера. Затем сломайте один тест намеренно, отправьте изменение в отдельной ветке и посмотрите на красный конвейер и предупреждение о непройденных проверках; после этого верните тест в рабочее состояние и убедитесь, что конвейер позеленел. Без аккаунта ту же пару состояний получите командой \texttt{mvn verify}.

\emph{Ответьте письменно:} почему сервер отклонил первую отправку и что сделала с вашим коммитом команда \texttt{git pull --rebase}; чем \texttt{git push --force-with-lease} лучше обычной принудительной отправки; зачем кешируется локальный репозиторий Maven; что произойдёт с последующими стадиями, если задание стадии \texttt{test} завершится с ошибкой.

\subsection*{Задание 5. Учебный кейс: Agile и Scrum}

Выберите одну тему из перечня курсовых работ: справочная система библиотеки, кинотеатра, магазина, ресторана, аэропорта или пиццерии. Ответы оформите в файле \texttt{docs/scrum-case.md} того же репозитория и закоммитьте его.

Начните с основы: выпишите четыре ценности манифеста в форме «А важнее Б», поясните оговорку «не отрицая важности того, что справа» и приведите ситуацию из вашего проекта, где правая часть перевешивает левую; опишите, как проект выглядел бы по каскадной модели и к какой потере это приведёт. Дальше опишите команду из пяти человек: кто владелец продукта, кто скрам-мастер, кто разработчики, и по две-три обязанности каждой роли \emph{в вашем проекте}; приведите пример решения владельца продукта и пример решения команды. Составьте бэклог продукта из восьми пользовательских историй по образцу листинга~\ref{lst:12-14} с оценкой по ряду Фибоначчи; для трёх верхних допишите критерии приёмки.

Спланируйте спринт: средняя скорость команды~--- 20~поинтов, спринт~--- две недели. Сформулируйте цель спринта одной фразой, выберите истории так, чтобы суммарная оценка не превышала 20~поинтов, и объясните выбор. Разбейте одну взятую историю на четыре-шесть технических задач так, как они будут выглядеть на доске, и укажите, что вы \emph{не} берёте в спринт и почему.

Завершите кейс организацией процесса. Составьте таблицу всех пяти событий для вашего проекта: событие, когда проходит, сколько длится, кто участвует, чем заканчивается. Сформулируйте отличие обзора от ретроспективы и напишите определение готовности из шести пунктов, включив в него ревью и зелёный конвейер из задания~4. Свяжите всё с Git: придумайте имена веток для трёх верхних историй, перечислите проверки, которые конвейер должен пройти до слияния, и напишите пример сообщения коммита, закрывающего задачу в трекере. Заведите доску задач со столбцами «Бэклог», «В работе», «Ревью», «Готово», перенесите туда бэклог спринта, проставьте столбцу «В работе» WIP-лимит, равный числу разработчиков, и опишите, что должно быть выполнено, чтобы карточка перешла из «В работе» в «Ревью».

\emph{Ответьте письменно:} почему подход называют набором ценностей, а не методологией, и чем тогда является Scrum; почему скрам-мастер не должен распределять задачи внутри команды; почему стори-поинты нельзя переводить в часы; какую проблему решает определение готовности.

\subsection*{Результаты занятия}

По итогам занятия у вас должны быть готовы: репозиторий \texttt{library} с историей не менее чем из двадцати коммитов, ветками, коммитом слияния, вручную разрешённым конфликтом и тегом \texttt{v1.0.0}; репозиторий \texttt{rebase-demo} с двумя ветками, слитой и перебазированной, и вывод \texttt{git log --oneline --graph --all} для обоих с пояснением, где видно слияние, а где перебазирование; журнал восстановления истории через reflog; файлы \texttt{pom.xml}, теста и конвейера с выводом успешной и падающей сборки; файл \texttt{docs/scrum-case.md} и снимок доски задач; ответы на вопросы всех пяти заданий.

\kontrolvoprosy

\begin{enumerate}
\item Чем распределённая система контроля версий отличается от
      централизованной? Назовите два следствия для повседневной работы.
\item Почему Git хранит снимки проекта, а не разницы между версиями?
      Перечислите четыре типа объектов и скажите, что хранит каждый.
\item Что такое ветка с точки зрения устройства репозитория, что такое HEAD и
      когда возникает состояние отсоединённого HEAD?
\item Что такое индекс и зачем он нужен, если можно было бы коммитить прямо из
      рабочего каталога?
\item Чем отличаются режимы \texttt{--soft}, \texttt{--mixed} и \texttt{--hard}
      команды \texttt{git reset}? Почему историю, которую видели другие,
      отменяют через \texttt{git revert}?
\item Сформулируйте золотое правило перебазирования и объясните, что именно
      ломается при его нарушении.
\item Что хранит reflog и когда он не поможет? Как \texttt{git bisect} находит
      виновный коммит примерно за восемь шагов?
\item Чем непрерывная интеграция отличается от непрерывной доставки, а
      доставка~--- от непрерывного развёртывания?
\item Назовите три роли Scrum и объясните, кто решает «что делать», а кто~---
      «как делать». Чем обзор спринта отличается от ретроспективы?
\item Назовите четыре ценности манифеста. Почему гибкий подход~--- это не отказ
      от документации и планирования?
\end{enumerate}

\testzadaniya

\begin{enumerate}

\item Чем распределённая система контроля версий принципиально отличается от
      централизованной?
  \begin{enumerate}
    \item[а)] она умеет хранить только текстовые файлы
    \item[б)] у каждого участника есть полная копия репозитория со всей историей
    \item[в)] изменения сохраняются автоматически, без явных коммитов
    \item[г)] история хранится только на сервере и раздаётся клиентам
  \end{enumerate}

\item В каком объекте Git хранится имя файла?
  \begin{enumerate}
    \item[а)] в blob, рядом с содержимым
    \item[б)] в объекте commit, вместе с сообщением
    \item[в)] в объекте tag
    \item[г)] в объекте tree
  \end{enumerate}

\item Какую роль играет индекс (staging area)?
  \begin{enumerate}
    \item[а)] ускоряет поиск по содержимому файлов
    \item[б)] хранит список всех веток и тегов
    \item[в)] кеширует объекты, скачанные с сервера
    \item[г)] промежуточная область, в которой собирается следующий коммит
  \end{enumerate}

\item Файл был закоммичен неделю назад, а сегодня вы добавили его имя в
      \texttt{.gitignore}. Что произойдёт?
  \begin{enumerate}
    \item[а)] файл будет удалён из репозитория при следующем коммите
    \item[б)] Git откажется коммитить, пока правило не убрано
    \item[в)] ничего: правила действуют только на неотслеживаемые файлы
    \item[г)] файл останется, но Git перестанет замечать его изменения
  \end{enumerate}

\item Вы создали ветку от основной, сделали два коммита, а основная за это
      время не менялась. Что произойдёт при слиянии?
  \begin{enumerate}
    \item[а)] Git создаст коммит слияния с двумя родителями
    \item[б)] Git откажется сливать и предложит перебазирование
    \item[в)] быстрая перемотка: указатель просто переместится вперёд
    \item[г)] возникнет конфликт, так как база слияния не определена
  \end{enumerate}

\item Что находится между маркерами \texttt{<<<<<<<} и \texttt{=======}?
  \begin{enumerate}
    \item[а)] версия из ветки, которую вы сливаете
    \item[б)] версия из текущей ветки, в которую выполняется слияние
    \item[в)] версия из общего предка обеих веток
    \item[г)] вариант объединения, предложенный Git автоматически
  \end{enumerate}

\item Что делает \texttt{git rebase main}, выполненная в ветке
      \texttt{feature}?
  \begin{enumerate}
    \item[а)] переносит коммиты \texttt{feature} поверх вершины \texttt{main},
      создавая новые коммиты с новыми хешами
    \item[б)] сливает \texttt{main} в \texttt{feature} коммитом слияния
    \item[в)] переносит указатель \texttt{main} на вершину \texttt{feature}
    \item[г)] удаляет из \texttt{feature} все коммиты, кроме последнего
  \end{enumerate}

\item Вы выполнили \texttt{git reset --hard HEAD\textasciitilde 2} и поняли,
      что потеряли два нужных коммита. Что поможет?
  \begin{enumerate}
    \item[а)] \texttt{git reflog}: найти прежнюю позицию HEAD и вернуть ветку
    \item[б)] \texttt{git revert} последнего коммита
    \item[в)] ничего: сброс удаляет коммиты безвозвратно
    \item[г)] получение изменений с сервера восстановит их в любом случае
  \end{enumerate}

\item Чем непрерывная доставка отличается от непрерывного развёртывания?
  \begin{enumerate}
    \item[а)] доставка собирает проект, а развёртывание прогоняет тесты
    \item[б)] доставка работает с тестовым стендом, развёртывание~--- с машиной
      разработчика
    \item[в)] доставка выполняется по расписанию, развёртывание~--- после
      каждого коммита
    \item[г)] при доставке решение о выкладке принимает человек, при
      развёртывании она происходит автоматически
  \end{enumerate}

\item Кто в Scrum определяет содержание и порядок бэклога продукта?
  \begin{enumerate}
    \item[а)] скрам-мастер
    \item[б)] команда разработчиков
    \item[в)] руководитель отдела разработки
    \item[г)] владелец продукта
  \end{enumerate}

\item В чём главное отличие Kanban от Scrum?
  \begin{enumerate}
    \item[а)] в Kanban задачи оцениваются в часах
    \item[б)] в Kanban нет спринтов и обязательных ролей: есть непрерывный
      поток задач и WIP-лимиты на столбцах доски
    \item[в)] в Kanban нет доски задач, работа ведётся по списку
    \item[г)] Kanban не подходит для разработки программного обеспечения
  \end{enumerate}

\end{enumerate}
